% ===================================================================
%  GAMBAR No. 5 — Omah lan Pambangunane.
%
%  Traduction, faite en 2026, en javanais d'aujourd'hui, sur l'ido de
%  texto/io. Regles de la colonne : voir 10-gambar-01.tex. Une seule
%  numerotation, et le plus long tableau du livret : 69 blocs, 142
%  renvois.
%
%  CE TABLEAU EST UN DIALOGUE D'UN BOUT A L'AUTRE, ET C'EST CE QUI
%  CHANGE TOUT POUR CETTE COLONNE-CI. Les niveaux du javanais ne se
%  choisissent pas par TEXTE mais par QUI PARLE A QUI. Tant que le
%  livret racontait, la regle « la colonne tient le ngoko » suffisait.
%  Ici Ioannes parle a son oncle, a monsieur Blanko et a l'architecte,
%  et un enfant javanais parle KRAMA a un adulte : il dit kula, mboten,
%  menika, sampun, matur nuwun. Les adultes, eux, lui repondent en
%  ngoko : aku, kowe, ora, iki, wis. L'asymetrie est la forme JUSTE, et
%  elle est meme ce que l'alinea decrit — l'oncle dit de l'enfant qu'il
%  est « tre serioza » et qu'il faut tout lui expliquer.
%
%  kolonoj.py A DONC ETE REPRIS AVANT D'ECRIRE CE TABLEAU. Les huit
%  regles de niveau — kula, mboten, menika, ingkang, sampun, kaliyan,
%  griya, toya — sont passees de « mot » a « narracio », une clef
%  nouvelle qui ne s'applique qu'aux fichiers NARRES. Le dialogue se
%  reconnait au fichier et la mesure est nette : les attributions de
%  parole s'ecrivent « \textsc{...}. --- », le tableau 5 en compte
%  TRENTE-SIX, le tableau 12 en compte UNE — et c'est « Noto. », pas un
%  locuteur —, les quatorze autres n'en ont aucune. Le seuil est pose a
%  cinq, ou il n'y a rien a departager.
%
%  ET LA REGLE N'A PAS ETE OTEE, ELLE A ETE DEPLACEE. Verifie sur le
%  tableau 4 en y glissant expres « kula mboten sampun ingkang » : les
%  quatre sont releves. Verifie sur celui-ci : aucun. C'est la
%  difference entre desarmer un controle et lui apprendre ou il
%  s'applique — le premier reflexe, au tableau 3, avait ete de
%  supprimer la regle « banyak », et c'etait le bon parce que le mot
%  n'etait pas le meme mot ; ici le mot EST le meme mot, et c'est la
%  situation qui change.
%
%  LE HOKKIEN PARAIT ICI, COMME IL AVAIT PARU AU MEME TABLEAU DANS LA
%  COLONNE INDONESIENNE. « cet » la peinture (103) vient de 漆, par les
%  peintres et les artisans chinois de Java — c'est exactement le mot
%  que l'en-tete du tableau 5 indonesien avait releve pour corriger son
%  compte de couches. Le tableau des metiers du batiment est celui qui
%  fait paraitre cette couche-la, dans les deux langues, et pour la
%  meme raison : ce sont les memes artisans.
%
%  MAIS C'EST LE NEERLANDAIS QUI TIENT LE CHANTIER. steger l'echafaud
%  (110), ember le seau (114), sekop la pelle (82), kwas le pinceau
%  (104), dhempul le mastic (73), gordhen le rideau, grendel le verrou
%  (53), seng le zinc (24), kakus les cabinets (20), balkon, garasi,
%  plester. Un batiment se fait avec les mots de celui qui l'a bati,
%  et a Java c'est le Hollandais qui a bati en dur.
%
%  ET POURTANT L'OUTIL DE LA MAIN RESTE JAVANAIS, ET LA REPARTITION EST
%  TRES REGULIERE : wadung la hache (116), pasah le rabot (91), tatah
%  le ciseau (94 bis), graji la scie (94), pacul la beche (131), garu le
%  rateau (130), gembor l'arrosoir (132), paron l'enclume (127), prapen
%  la forge (126), sapit la tenaille (128), cetok la truelle. Le
%  neerlandais a donne l'ECHAFAUDAGE et le batiment ; le javanais a
%  garde ce qui tient dans la main. La meme frontiere separe wedhi le
%  sable (62) et gamping la chaux (63), qui sont d'ici, de plester et
%  bata, qui sont venus.
%
%  ET LE TOIT EST ENTIEREMENT JAVANAIS : payon le toit (31), genthing
%  la tuile (55), pyan le plafond (26), pandhe le forgeron, pawon la
%  cuisine (d), latar la cour (45), undhak-undhakan les marches (14),
%  tlundhak le seuil (10). L'indonesien dit atap, genting, langit-langit,
%  dapur, halaman — et sur les huit, deux seulement se ressemblent.
% ===================================================================

%%K t05-apar-1 apar
\VUsaut{6.0mm}
\VUtitre{15.0pt}{60}{GAMBAR No. 5}
\VUsaut{1.4mm}
\VUtitre{11.4pt}{0}{Omah lan Pambangunane.}
\VUsaut{2.2mm}

%%K t05-tit-1 sub
\VUpk{10.2pt}{40}{Kepethuk kang becik.}
\VUsaut{3.0mm}

%%K t05-01-1 p
\textsc{Ioannes}. --- Pakdhe, kula kepengin sanget mangertos kepripun
tiyang mbangun \VUgras{omah}.

%%K t05-01-2 p
\textsc{Pakdhe} (80). --- Iku gampang banget, kancaku ; ayo melu aku,
dakduduhi omah kang lagi dibangun dening Pak Bernardus (79) lan kang
bakal dakenggoni.

%%K t05-01-3 p
\textsc{Ioannes}. --- Sinten ingkang nggambar \VUgras{denah}
\textsuperscript{(76)} omah menika?

%%K t05-01-4 p
\textsc{Pakdhe}. --- \VUgras{Arsitek} \textsuperscript{(75)} ; dheweke
masrahake gambar kang cetha banget lan gambar iku disarujuki, banjur
\VUgras{pemborong} \textsuperscript{(77)} miwiti pagaweane.

%%K t05-01-5 p
\textsc{Ioannes}. --- Lajeng sinten pemborong menika?

%%K t05-01-6 p
\textsc{Pakdhe}. --- Yaiku Pak Blanko, bapakne kancamu Iakobus. Dheweke
mimpin para tukang bareng karo Pak Rufo, sang arsitek.

%%K t05-01-7 p
Lah! Iki Pak Blanko teka pas ing wektune karo \VUgras{garisane}
\textsuperscript{(78)}, lan Pak Rufo karo denahe.

%%K t05-01-8 p
Sugeng enjing, Bapak-bapak. Kepripun kabare?

%%K t05-01-9 p
\textsc{Pak Rufo}. --- Sae sanget, matur nuwun. Sugeng enjing,
Ioannes ; wis gedhe temen bocah iki saiki!

%%K t05-01-10 p
\textsc{Pakdhe}. --- Iya, tenan, wis kaya wong lanang cilik. Dheweke
tenanan banget ; kabeh kang dideleng kudu diterangake marang dheweke.
Dina iki dheweke kepengin ngerti kepriye wong mbangun omah.

%%K t05-tit-2 sub
\VUpk{10.2pt}{40}{Para tukang.}
\VUsaut{3.0mm}

%%K t05-01-11 p
\textsc{Pak Blanko}. --- Lah, melua aku, kanca cilik ; sedhela maneh
dakterangake marang kowe, amarga aku seneng banget marang bocah kang
kepengin sinau.

%%K t05-01-12 p
Kowe weruh tukang kang nyekel \VUgras{sekop} \textsuperscript{(82)} ing
tangane iku : dheweke \VUgras{tukang dhudhuk} \textsuperscript{(81)}.
Dheweke ndhudhuk \VUgras{kalen} \textsuperscript{(46)} kanggo masang pipa
banyu. Dheweke uga ndhudhuk lemah ing panggonan omah iku ngadeg,
supaya tukang batu bisa ngadegake \VUgras{tembok} papat. Perangan
tembok kang ana ing jero lemah diarani \VUgras{pondhasi}
\textsuperscript{(2)}. Ruwang kang kinurung ing antarane pondhasi iku
dadi \VUgras{kamar ngisor lemah} \textsuperscript{(3)}. Kamar ngisor
lemah iku duwe jendhela cilik kang diarani \VUgras{jendhela kamar
ngisor lemah} \textsuperscript{(5)}, \VUgras{lawang}
\textsuperscript{(4)} lan andha.

%%K t05-01-13 p
\VUgras{Tukang batu} \textsuperscript{(108)} nerusake mbangun tembok iku
karo ninggal bolongan kanggo \VUgras{lawang} \textsuperscript{(8)} lan
\VUgras{jendhela} \textsuperscript{(17)}.

%%K t05-01-14 p
\textsc{Ioannes}. --- Nanging tukang batu menika mboten kula tingali!

%%K t05-01-15 p
P\textsuperscript{k} \textsc{Blanko}. --- Saiki dheweke wis rampung
nyambut gawe ing omah iku. Dheweke ana ing cedhake \VUgras{kandhang
jaran} \textsuperscript{(54)}, ing dhuwur \VUgras{stegere}
\textsuperscript{(110)} ; dheweke mbangun temboke \VUgras{papan
umbah-umbah} \textsuperscript{(56)}.

%%K t05-01-16 p
Dheweke duwe cetok, bak luluh lan \VUgras{unting-unting}
\textsuperscript{(109)}. Nanging jeneng-jeneng iku ora bakal kokeling.

%%K t05-01-17 p
\textsc{Ioannes}. --- Mesthi kemawon kula enget, amargi kula badhe
nyuwun dhateng pakdhe kula cetok alit saha bak, lan unting-unting badhe
kula damel piyambak mawi tali.

%%K t05-tit-3 sub
\VUsaut{3.0mm}
\VUpk{10.2pt}{40}{Bahan bangunan.}
\VUsaut{3.0mm}

%%K t05-01-18 p
\textsc{Pakdhe}. --- Ha! becik temen. Nanging kandhaa marang aku,
Ioannes, apa kang dibutuhake kanggo mbangun omah?

%%K t05-01-19 p
\textsc{Ioannes}. --- Ingkang dipunbetahaken \VUgras{watu tembok}
\textsuperscript{(59)} saha \VUgras{luluh} \textsuperscript{(65)}.

%%K t05-01-20 p
\textsc{Pakdhe}. --- Bener banget ; delengen wong kang nganggo topi
gedhe ing dhuwur \VUgras{grobage} \textsuperscript{(136)} iku. Dheweke
\VUgras{tukang grobag} \textsuperscript{(135)}. Saben dina dheweke
nggawa \VUgras{bongkahan watu} \textsuperscript{(60)} mrene. Dheweke
ngedhunake momotane ing latar, lan \VUgras{tukang natah watu}
\textsuperscript{(83)} natahi bongkah iku lan nggraji nganggo
\VUgras{graji watu} \textsuperscript{(85)}. \VUgras{Kuli}
\textsuperscript{(146)} nggawa watu iku menyang tukang batu, kang
masang.

%%K t05-01-21 p
Sauntara iku, \VUgras{tukang ngaduk luluh} \textsuperscript{(87)}
nyiyapake luluh. Dheweke butuh \VUgras{wedhi} \textsuperscript{(62)},
\VUgras{gamping} \textsuperscript{(63)} lan \VUgras{banyu}
\textsuperscript{(64)}. Gamping ana ing cedhake ing \VUgras{karung}
\textsuperscript{(71)} ; \VUgras{laden tukang batu}
\textsuperscript{(111)} nggawakake wedhi marang dheweke ing
\VUgras{grobag dorong} \textsuperscript{(112)}, lan \VUgras{bocah
magang} \textsuperscript{(113)} ngadeg ing pompa (58). Bocah iku ngiseni
\VUgras{ember} \textsuperscript{(114)}. Luluh iku sedhela maneh wis
siyap. \VUgras{Tukang usung luluh} \textsuperscript{(107)} njupuk lan
nggawa munggah, liwat andha, menyang dhuwur steger, panggonan tukang
batu ngrampungake sisihe omah iku.

%%K t05-01-22 p
Lan saiki, apa jenenge tukang kang nggarap \VUgras{payon}
\textsuperscript{(31)}?

%%K t05-01-23 p
\textsc{Ioannes}. --- Piyambakipun \VUgras{tukang payon}
\textsuperscript{(118)}.

%%K t05-tit-4 sub
\VUsaut{3.0mm}
\VUpk{10.2pt}{40}{Payone omah.}
\VUsaut{3.0mm}

%%K t05-01-24 p
P\textsuperscript{k} \textsc{Blanko}. --- Iya, nanging kang teka
ndhisik yaiku \VUgras{tukang kayu} \textsuperscript{(115)}.

%%K t05-01-25 p
Tukang kayu nggarap \VUgras{rangka kayu} \textsuperscript{(35)}.
Dheweke nyambung \VUgras{balok} \textsuperscript{(37)} nganggo
\VUgras{pathok} \textsuperscript{(117)}. Para tukang ing ndhuwur kono,
kang nganggo kathok beludru amba, iku tukang kayu. Kang siji nyekel
balok cilik, sijine ngethok pathok nganggo \VUgras{wadunge}
\textsuperscript{(116)}. Kang ana ing cedhake \VUgras{cerobong asap}
\textsuperscript{(41)} iku tukang payon. Dheweke mathok reng ing dhuwur
(*) balok, lan \VUgras{watu tulis} \textsuperscript{(66)} ing dhuwur
reng. Kadhang dheweke masang genthing.

%%K t05-noto-1 noto
\VUnotes{143.2mm}{%
(*) \VUgras{Balk-o} = J. \textit{Balken}, I. \textit{joist}, P.
\textit{solive}, It. \textit{travicello}, R. \textit{balka}, Sp. Port.
\textit{viga}.
Oyod teknis kang nganti saiki durung ditampa, nanging pantes
diusulake kanggo negesake tembung P. \textit{solive}, kang dudu trabo
(P. \textit{poutre}) lan dudu uga trabeto. Iku watang kayu pesagi kang
nyangga jogan papan lan pyan.%
}

%%K t05-01-26 p
Payone kandhang jaran iku \VUgras{genthingan} \textsuperscript{(55)},
payone omah \VUgras{watu-tulisan} \textsuperscript{(32)}, lan payone
emper saka \VUgras{seng} \textsuperscript{(24)}.

%%K t05-01-27 p
\textsc{Ioannes}. --- Punapa tukang payon ugi nggarap payon seng?

%%K t05-01-28 p
P\textsuperscript{k} \textsc{Blanko}. --- Ora ; kang nggarap yaiku
\VUgras{tukang seng} \textsuperscript{(121)} utawa \VUgras{tukang
timbel,} yaiku kang masang \VUgras{jendhela payon}
\textsuperscript{(38)} ing dhuwure payon omah. Dheweke uga masang
\VUgras{talang} \textsuperscript{(43)} lan \VUgras{talang payon}
\textsuperscript{(42)}. Dheweke nyolder seng lan pelat. Dheweke kang
masang \VUgras{panulak bledheg} \textsuperscript{(40)} lan
\VUgras{panuduh arah angin} \textsuperscript{(39)} ing dhuwur
\VUgras{tembok pucuk} \textsuperscript{(36)}.

%%K t05-01-29 p
\textsc{Ioannes}. --- Dados omahipun sampun rampung?

%%K t05-tit-5 sub
\VUsaut{3.0mm}
\VUpk{10.2pt}{40}{Piranti.}
\VUsaut{3.0mm}

%%K t05-01-30 p
P\textsuperscript{k} \textsc{Blanko}. --- Durung kabeh. \VUgras{Tukang
plester} \textsuperscript{(99)} mbangun nganggo \VUgras{bata}
\textsuperscript{(61)} tembok-tembok kang mbagi omah iku dadi
warna-warna kamar. Nganggo \VUgras{plester} \textsuperscript{(70)}
dheweke nglapisi \VUgras{pyan} \textsuperscript{(26)} lan tembok. Saiki
dheweke nggarap pyane \VUgras{emper} \textsuperscript{(33)}. Kaya tukang
batu, dheweke duwe \VUgras{cetok} \textsuperscript{(100)} lan \VUgras{bak
luluh} \textsuperscript{(101)}.

%%K t05-01-31 p
Banjur tukang kayu alus nggarap lawang, jendhela, \VUgras{jogan papan}
\textsuperscript{(16)} lan \VUgras{tutup jendhela} \textsuperscript{(19)}.

%%K t05-01-32 p
Saiki dheweke nyambut gawe ing latar ing dhuwur \VUgras{meja gawene}
\textsuperscript{(90)}. Dheweke duwe \VUgras{graji}
\textsuperscript{(94)} kanggo nggraji kayu (69), \VUgras{pasah}
\textsuperscript{(91)}, \VUgras{penjepit} \textsuperscript{(92)},
\VUgras{tatah} \textsuperscript{(94 bis)}, \VUgras{jangka}
\textsuperscript{(95)} lan \VUgras{palu} kayu \textsuperscript{(95 bis)}.

%%K t05-01-33 p
\textsc{Ioannes}. --- Ha! jebul langkung nyenengaken nggarap kayu
tinimbang masang watu. Pakdhe, tumbasaken kula meja gawe, graji saha
pasah, amargi kula badhe damel kandhang kangge Medor.

%%K t05-01-34 p
\textsc{Pakdhe}. --- Iya, le.

%%K t05-01-35 p
\textsc{Ioannes}. --- Bingah sanget kula! Lajeng sinten ingkang ngadeg
ing cakete lawang ngajeng menika?

%%K t05-tit-6 sub
\VUsaut{3.0mm}
\VUpk{10.2pt}{40}{Ngecet.}
\VUsaut{3.0mm}

%%K t05-01-36 p
P\textsuperscript{k} \textsc{Blanko}. --- Dheweke \VUgras{tukang cet}
\textsuperscript{(102)}. Dheweke ngadeg ing dhuwur andhane karo
sakaleng \VUgras{cet} \textsuperscript{(103)} lan \VUgras{kwase}
\textsuperscript{(104)}. Dheweke ngecet kayune lawang ngarep supaya
luwih awet.

%%K t05-01-37 p
Dheweke uga nemplekake kertas tembok ing kamar-kamar.

%%K t05-01-38 p
\VUgras{Tukang gordhen} \textsuperscript{(107)} kang ana ing dhuwur
\VUgras{andha} \textsuperscript{(106)} ing \VUgras{balkon}
\textsuperscript{(28)} lagi masang \VUgras{tirai} \textsuperscript{(29)}.

%%K t05-01-39 p
Tukang kang ngadeg ing cedhake jendhela gedhe iku \VUgras{tukang kaca}
\textsuperscript{(96)}. Dheweke masang \VUgras{kaca jendhela}
\textsuperscript{(18)} nganggo \VUgras{dhempul} \textsuperscript{(73)}.
Tukang sijine, kang jengkeng ing cedhake lawang kamar ngisor lemah,
yaiku \VUgras{tukang kunci} \textsuperscript{(97)}. Dheweke masang
\VUgras{kunci} \textsuperscript{(6)}. \VUgras{Kikire}
\textsuperscript{(98)} ana ing cedhake.

%%K t05-01-40 p
\textsc{Ioannes}. --- Punapa tukang ingkang nggebug wesi mawi
\VUgras{palunipun} \textsuperscript{(84)} menika ugi tukang kunci?

%%K t05-01-41 p
P\textsuperscript{k} \textsc{Blanko}. --- Luwih trepe, dheweke pandhe
wesi (125). Dheweke nempa \VUgras{barang wesi} \textsuperscript{(67)}
kanggo \VUgras{tukang listrik} \textsuperscript{(124)} kang ana ing
dhuwure payon \VUgras{garasi} \textsuperscript{(51)}. Dheweke duwe
\VUgras{prapen} \textsuperscript{(126)}, \VUgras{paron}
\textsuperscript{(127)} lan \VUgras{sapit} \textsuperscript{(128)}.

%%K t05-01-42 p
Tukang listrik masang \VUgras{kawat} \VUgras{tembaga}
\textsuperscript{(44)} telpon.

%%K t05-tit-7 sub
\VUpk{10.2pt}{40}{Nggarap kebon.}
\VUsaut{3.0mm}

%%K t05-01-43 p
\textsc{Pakdhe}. --- Ing pojoke \VUgras{latar} \textsuperscript{(45)},
ing cedhake \VUgras{pager wesi} \textsuperscript{(47)} lan
\VUgras{gapura} \textsuperscript{(48)}, kowe weruh \VUgras{tukang
kebon} \textsuperscript{(129)}. Dheweke nyiyapake \VUgras{kebon cilik}
\textsuperscript{(49)}. Dheweke wis macul lemah nganggo \VUgras{pacule}
\textsuperscript{(131)}, dheweke nyebar wiji lan nggaru nganggo
\VUgras{garune} \textsuperscript{(130)}. Banjur, nganggo \VUgras{gembore}
\textsuperscript{(132)}, dheweke bakal ngeleb \VUgras{bedhengan}
\textsuperscript{(150)} lan tanduran iku bakal thukul.

%%K t05-01-44 p
\VUgras{Tukang jaran} \textsuperscript{(133)}, kang lagi wae ngopeni
jaran lan menehi pakan, ngresiki \VUgras{kreta} \textsuperscript{(52)}
nganggo spons, \VUgras{pompa tangan} \textsuperscript{(134)} lan
saember banyu. Banjur dheweke bakal nglebokake kreta iku maneh menyang
garasi lan nutup lawange nganggo \VUgras{grendel} \textsuperscript{(53)}
kang kandel.

%%K t05-01-45 p
Lah, dakkira kowe wis marem ; kowe wis oleh katrangan kang cukup.

%%K t05-01-46 p
\textsc{Ioannes}. --- Inggih, Pakdhe, nanging kula kepengin sanget
ningali njeronipun omah menika.

%%K t05-01-47 p
\textsc{Pakdhe}. --- Iku bakal kelakon sesuk. Pak Rufo bakal
nerangake denahe marang kowe lan nuduhake jenenge saben kamar.

%%K t05-tit-8 sub
\VUsaut{3.0mm}
\VUpk{10.2pt}{40}{Tata ruwang jero omah.}
\VUsaut{2.4mm}

%%K t05-apar-2 apar
{\centering\textit{(Delengen denahe.)}\par}
\VUsaut{2.4mm}

%%K t05-01-48 p
\textsc{Arsitek}. --- Ayo, Ioannes, sedhela maneh dakajak kowe
njlajahi perangan-perangane omah iki.

%%K t05-01-49 p
Ayo padha mlebu \VUgras{lantai ngisor} \textsuperscript{(7)}. Iki
tombole \VUgras{bel} \textsuperscript{(11)}. Kita munggah telung
\VUgras{undhak-undhakan} \textsuperscript{(14)} \VUgras{undhakan
ngarep} \textsuperscript{(9)}, kita ngliwati \VUgras{tlundhak}
\textsuperscript{(10)} \VUgras{lawang ngarep} \textsuperscript{(8)}, lan
iki kita ana ing sikile \VUgras{andha} \textsuperscript{(13)}.

%%K t05-01-50 p
\textsc{Ioannes}. --- Menika lorongipun.

%%K t05-01-51 p
\textsc{Arsitek}. --- Dudu, iki \VUgras{ruwang ngarep}
\textsuperscript{(12)}. \VUgras{Lorong} \textsuperscript{(a)} ana ing
pucuk kana. Ing sisih tengen, iki \VUgras{ruwang tamu}
\textsuperscript{(b)} lan \VUgras{ruwang mangan} \textsuperscript{(c)} ;
adhep-adhepan karo ruwang mangan ana \VUgras{pawon}
\textsuperscript{(d)} lan \VUgras{kamar batur} \textsuperscript{(e)},
banjur ing ngarepe \VUgras{kamar nulis} \textsuperscript{(f)}.
\VUgras{Emper} \textsuperscript{(23)} karo \VUgras{lawang kacane}
\textsuperscript{(25)} ana ing cedhake ruwang tamu.

%%K t05-01-52 p
Saiki kita bakal munggah andha. Cekelan ing \VUgras{cekelan andha}
\textsuperscript{(15)} lan etungen undhak-undhakane. Ana patbelas. Iki
\VUgras{selasar} \textsuperscript{(m)} ; kita ana ing \VUgras{lantai}
\textsuperscript{(27)} kapisan.

%%K t05-tit-9 sub
\VUsaut{3.0mm}
\VUpk{10.2pt}{40}{Lantai kapisan.}
\VUsaut{3.0mm}

%%K t05-01-53 p
Adhep-adhepan karo andha ana \VUgras{kakus} \textsuperscript{(20)} lan
\VUgras{kamar rias} \textsuperscript{(j)}. Ing sisih tengen ana
\VUgras{kamar turu} \textsuperscript{(g)} kang endah karo jendhela loro
kang madhep \VUgras{ngarepe omah} \textsuperscript{(1)}. Ing sisih
kiwa ana \VUgras{kamar adus} \textsuperscript{(1)} lan \VUgras{kamar
tamu} \textsuperscript{(i)} karo \VUgras{ceruk turu}
\textsuperscript{(h)} kang gedhe. Ing cedhake kamar turu ana
\VUgras{kamar bocah} \textsuperscript{(k)}. Kamar iku madhep
\VUgras{teras} \textsuperscript{(21)} kang jembar. Ing sangisore teras
iku ana kakus sijine (20), lan ing tembok, iki \VUgras{lonceng}
\textsuperscript{(22)} kang muni kanggo bujana bengi.

%%K t05-01-54 p
\textsc{Ioannes}. --- Sumangga kita minggah dhateng \VUgras{lantai
kalih}.

%%K t05-01-55 p
\textsc{Arsitek}. --- Ora ana lantai loro. Ing sangisore payon mung
ana \VUgras{kamar loteng} \textsuperscript{(33)} lan loteng (34).
Panjlajahan kita wis rampung ; kita kena mudhun maneh.

%%K t05-01-56 p
\textsc{Ioannes}. --- Matur nuwun, Pak ; sedaya menika narik sanget.
Kula badhe enggal nyariyosaken dhateng bapak kula sedaya ingkang kula
tingali.

\VUsaut{4.0mm}
\VUfilet{22mm}
